\section{Nilpotent blocks and a common projector}\label{sec:block-algebra}
The transfer proof reduces the whole theorem to Exact Deletion
(Theorem~\ref{thm:deletion}). We now prove that finite-dimensional input,
using a support budget $r$ as an abstract parameter. For general indexed
targets its application is $r=2kd+s+t$; for invariant targets it is
$r=kd+s$. The construction is independent of the equations and targets.
We begin with the simpler square-zero operator in characteristic two,
then build a witness with a quadratic retained-label cost, replace that
cost by the linear threshold, and finally handle odd characteristic.

There are two obligations. On permutation bases, deletion must retain
precisely the objects using the $n$ retained labels and transport the
actual coefficient maps. On their subspaces and quotients it must also
preserve kernels and images, so that it transports PC intersections
after cancellation. The orbit calculation supplies the first
obligation; one common group-algebra witness will supply the second.

\subsection{A block of labels and a square-zero operator}\label{sec:block2}
For this part $\operatorname{char}K=2$. Choose $q=2^a>r$, put $N=n+q$, and split the labels into retained labels $R=[n]$ and a block $B$ with $q$ labels. Identify $B$ with the vector space $A=\Ftwo^a$. Translation by the $i$th unit vector gives a permutation $\sigma_i$ of $B$, extended by the identity on $R$. These permutations commute and square to the identity.

Choose $c_1,\ldots,c_a$ linearly independent over $\Ftwo$ in a suitable extension field, and put
\begin{equation}\label{eq:z2}z=\sum_{i=1}^ac_i(\sigma_i-1).\end{equation}
Here $z$ is a linear operator, not a scalar, and its products mean composition. Thus $z^2=0$ can hold although $z\ne0$. Since the summands commute, cross terms cancel in characteristic two, and
\[z^2=\sum_i c_i^2(\sigma_i-1)^2=0.\]
For $q=2$, $z$ can be the swap of the two block labels plus the identity. For $q=4$, work over $\mathbb F_4=\Ftwo[\omega]/(\omega^2+\omega+1)$ and use $(c_1,c_2)=(1,\omega)$.
During this finite-algebra construction, $K$ denotes the chosen
working extension field in all tuple spaces and group algebras.
The extension field is used to construct the block operator and its
witness; Lemma~\ref{lem:field} preserves derivability and minimum
degree over the original coefficient field.

The independence is essential. On a two-point orbit where $\sigma_1$ and $\sigma_2$ induce the same swap, choosing $c_1=c_2=1$ makes $z=0$. The orbit is moved, yet would survive deletion. Independent coefficients ensure that every nonzero binary combination of the $c_i$ remains nonzero.

For a space $M$ with label action define
\begin{equation}\label{eq:D2}\D(M)=\ker z/\im z.\end{equation}
This quotient is legal because $z^2=0$ implies $\im z\subseteq\ker z$. Two vectors killed by $z$ represent the same class if their difference is $zw$ for some $w$. A map commuting with label permutations commutes with $z$ and induces a map between these quotients. Permutations of the retained labels commute with $z$, so they still act on the quotient.

\subsubsection{The two-point orbit calculation}
For the illustrative choice $z=\sigma-I$, on basis vectors $e_\beta,e_\gamma$ swapped by $\sigma$,
\[ze_\beta=ze_\gamma=e_\beta+e_\gamma.\]
The kernel and image are both the one-dimensional span of $e_\beta+e_\gamma$, so the quotient is zero. On a fixed basis vector $e_i$, $z=0$, and its line survives. This is the prototype of deletion: moved directions disappear, fixed directions remain.

\subsubsection{The calculation for every moved orbit}
An orbit comprises the objects reached from one object by the block
translations. Its stabiliser consists of translations fixing that
object. Since $A$ is abelian, an orbit with stabiliser $H$ can be
coordinatised by the cosets $A/H$; the coordinate function below uses
that description.

Let $X$ be a finite $A$-set with no point fixed by all of $A$. On each orbit choose $x_0$. Its stabiliser $H$ is a proper binary subspace. Choose a nonzero linear functional $\ell:A\to\Ftwo$ vanishing on $H$. The function
\[\pi(vx_0)=\ell(v)\]
is well defined: two choices of $v$ differ by an element of $H$. Let $D$ multiply $e_x$ by $\pi(x)$, and put $e_i=\ell(\varepsilon_i)$. These are constants on that orbit. If $S_i$ denotes the action of $\sigma_i$, direct evaluation at a basis vector gives
\[DS_i-S_iD=e_iS_i,\qquad Dz-zD=W:=\sum_i c_ie_iS_i.\]
The difference $DS_i-S_iD$ is a \emph{commutator}, measuring the
failure of the two operators to commute. Here $D$ is diagonal in the
orbit basis, and $e_i\in\Ftwo$ is a scalar, distinct from a basis
vector $e_x$.

Moreover $Wz=zW$ and
\[W^2=\Bigl(\sum_i c_ie_i\Bigr)^2I=\alpha^2I,\qquad\alpha\ne0.\]
Thus $\Theta=W^{-1}D$ satisfies $\Theta z-z\Theta=I$. For $zm=0$ this gives $m=-z\Theta m$, hence $\ker z\subseteq\im z$. The reverse inclusion was automatic, so the quotient vanishes.

This calculation also works on $K[X]\otimes Y$ for \emph{any} space $Y$ carrying the label action. Let $D$ inspect only the $X$ coordinate and let $S_i$ act on both factors. We still have $S_i^2=I$ and the same commutator calculation. This stronger statement will justify multiplication, whose domain contains a variable factor tensored with a consequence subspace.

\subsection{Natural deletion on permutation bases}
Let $\mathcal U$ be a basis system whose objects have support less than $q$. An object fixed by all of $A$ cannot have support meeting $B$. Otherwise invariance of its support under the regular action would force its support to contain all $q$ block labels. Thus its fixed objects are precisely $\mathcal U_n$.

Split the basis into fixed and moved objects. The preceding calculation kills the moved summand and fixes the retained one. Hence
\begin{align}
\D(K[\mathcal U_{n+q}])&\cong K[\mathcal U_n],\label{eq:basis}\\
\D(K[\mathcal U_{n+q}]\otimes Y)&\cong K[\mathcal U_n]\otimes\D(Y).\label{eq:tensor}
\end{align}
The second identification sends $e_u\otimes[y]$ to $[e_u\otimes y]$ for retained $u$. It is natural in $Y$: applying a map to $y$ before or after taking its class gives the same result. For a product of two permutation bases, it simply keeps pairs of retained objects. There is no choice of unrelated isomorphisms here.

Suppose a matrix from one permutation basis to another has stable relabelling-invariant coefficients. On a retained input vector, its image is killed by $z$. Its retained component has exactly the smaller-rank coefficients; its moved component is in $\ker z=\im z$ on the moved summand, so vanishes in the quotient. Thus the \emph{actual matrix}, not merely its rank, is transported correctly.

\subsubsection{A small failure of exactness}
Take the two-point permutation space $E=Ke_\beta\oplus Ke_\gamma$ and its invariant line $T=K(e_\beta+e_\gamma)$. Inclusion $T\hookrightarrow E$ is injective. But $\D(T)=T$ whereas $\D(E)=0$. Deletion takes an injection to the zero map on a nonzero space. Therefore the basis calculation alone does not prove that deletion preserves kernels, images or intersections.

This failure is precisely why a proof cannot stop at ``the monomials delete correctly.'' We must show that, with enough retained labels, every subspace produced by the PC calculation lies in a class where the quotient operation is exact.

\subsection{A common group-algebra witness}\label{sec:common}
The group algebra $K[S_N]$ consists of finite formal sums $\sum_ga_gg$ of label permutations. It acts on every space with an $S_N$ action by the corresponding linear combination of operators. Its multiplication is composition. This provides a common language for certificates valid on many spaces at once.

For $b\ge0$ let
\[P_b(N)=K[\Inj([b],[N])],\]
the space on ordered injective $b$-tuples. This is a tuple space, distinct from the polynomial ambient $P_{N,\le d}$; $b$ counts positions rather than polynomial degree. Label permutations act on entries; a second group $S_b$ permutes tuple positions. These actions commute. Define
\[J_r(N)=\{a\in K[S_N]:aP_b(N)=0\text{ for every }b\le r\}.\]
It is a two-sided ideal. Write $\mathcal C_r(N)$ for the finite-dimensional spaces with label action annihilated by this ideal. Subspaces invariant under labels, quotients and finite direct sums remain in this class.
The ideal $J_r(N)$ is unrelated to the prescribed-target map $J_N$.

A two-sided ideal is closed under multiplication on either side by
any group-algebra element. Here this follows because an operator
annihilating every tuple space still does so after either composition.
Its relations therefore hold on the subspaces and quotients in
$\mathcal C_r(N)$.

Every bounded-support permutation basis of support at most $r$ belongs to the class. For an orbit of an object, order its support of size $b\le r$. Sending an injective tuple to the corresponding relabelling of that object gives a surjective equivariant map from $P_b(N)$ onto the orbit space. The support condition makes this well defined. Elements of $J_r$ kill that quotient as well.

\begin{definition}
A support-$r$ witness is one element $w\in K[S_N]$ satisfying
\begin{equation}\label{eq:witness2}zwz-z\in J_r(N).
\end{equation}
\end{definition}
Such a $w$ gives $zwz=z$ on \emph{every} object of $\mathcal C_r(N)$. It is a inner inverse: for $v\in\im z$, applying $w$ produces a preimage under $z$. It does not invert the nilpotent operator on the whole space. Choosing a separate matrix inverse on each space would not give this compatibility. A module map $f$, meaning a linear map commuting with label permutations, satisfies $f(wv)=w f(v)$ for this common group-algebra element.

\subsubsection{The projector calculation, without an exactness assumption}
Recall that a projector is an idempotent linear operator. Its image
can model a quotient without being invariant under the original label
action. What matters for exactness here is that the same formula for
$e$ commutes with every equivariant map between the spaces concerned.

Suppose $z^2=0$ and $zwz=z$. Set
\[e=(I-zw)(I-wz).\]
Then
\[ze=z(I-wz)=z-zwz=0,\qquad ez=(I-zw)z=0.\]
If $zv=0$, then $ev=v-zwv$. Also
\[(zw)e=(zw-zwzw)(I-wz)=0,\]
since $zwzw=zw$. Applying the formula on $\ker z$ to $ev$ gives $e^2v=ev$. Thus $e$ is a projector.

The map $[v]\mapsto ev$ is an isomorphism $\ker z/\im z\to eM$: it is well defined because $ez=0$, its image lies in $\ker z$, and $ev-v=-zwv$ is in $\im z$ when $zv=0$. Its inverse sends a vector of $eM$ to its quotient class.

A short exact sequence $0\to U\to V\to W\to0$ consists of an
injection, followed by a surjection whose kernel is the image of that
injection. Preservation means that these three properties still hold
after applying deletion. The common projector verifies each property
directly.

Now let $0\to U\to V\to W\to0$ be an exact sequence in $\mathcal C_r(N)$. All maps commute with the same $e$. Restricting to $eU,eV,eW$ preserves injectivity. A vector of $eW$ has a lift $v\in V$; $ev$ is a lift in $eV$, proving surjectivity. A vector $ev$ mapping to zero lies in $U$ and equals its own projection, so lies in $eU$. This proves exactness explicitly. By the natural isomorphism, $\D$ is exact.

The projector need not commute with permutations of the retained labels. We use its image as a vector-space model to prove exactness, and keep the original quotient for the residual label action. Nothing here asserts an unwarranted equivariant direct-sum decomposition for that residual group.

\begin{example}[One retained label]
Let the retained label be $a$ and the two block labels be $b,c$. Put $z=I+(bc)$ and $w=(ab)$. On $P_1(3)$,
\[zwz(e_b)=zw(e_b+e_c)=z(e_a+e_c)=e_b+e_c=ze_b,\]
and similarly for $e_c$; both sides vanish on $e_a$. The projector is the all-ones $3\times3$ matrix. Its image is the constant line, naturally representing the retained one-dimensional quotient. With no retained label, every element of $K[S_2]$ satisfies $zwz=0$, so there is no such witness. The spare label makes a real difference.
\end{example}
